\documentclass{ximera}

\input{../../../preamble.tex}


\definecolor{fvdnavy}{HTML}{071D43}
\definecolor{fvdcyan}{HTML}{20BFC6}
\definecolor{fvdcyanlight}{HTML}{EFFCFB}
\definecolor{fvdmagenta}{HTML}{F10D69}
\definecolor{fvdmagentalight}{HTML}{FFF1F7}
\definecolor{fvdtext}{HTML}{163044}





%=================================================
% Pedagogische omgevingen
% PDF: tcolorbox
% HTML: learning-box
%=================================================

\newenvironment{voorkennis}
{%
    \ifdefined\HCode
        \HCode{%
<section class="learning-box learning-box--prior">
<div class="learning-box__header">
<span class="learning-box__icon" aria-hidden="true"></span>
<span class="learning-box__title">Voorkennis</span>
</div>
<div class="learning-box__content">
        }%
        \begin{itemize}
    \else
       \begin{tcolorbox}[
    enhanced,
    breakable,
    colback=fvdcyanlight,
    colframe=fvdcyan,
    coltext=fvdtext,
    boxrule=0.5pt,
    leftrule=4pt,
    arc=4mm,
    outer arc=4mm,
    left=7mm,
    right=6mm,
    top=5mm,
    bottom=5mm,
    before skip=6mm,
    after skip=6mm,
    title=Voorkennis,
    coltitle=fvdcyan,
    fonttitle=\bfseries\Large,
    attach title to upper={\par\smallskip},
    boxed title style={
        empty,
        left=0mm,
        right=0mm,
        top=0mm,
        bottom=0mm
    },
    overlay={
        \draw[
            fvdcyan,
            line width=0.5pt
        ]
        ([xshift=42mm,yshift=-13mm]frame.north west)
        --
        ([xshift=-7mm,yshift=-13mm]frame.north east);
    }
]
        \begin{itemize}
    \fi
}
{%
    \end{itemize}
    \ifdefined\HCode
        \HCode{%
</div>
</section>
        }%
    \else
        \end{tcolorbox}
    \fi
}


\newenvironment{leerdoelen}
{%
    \ifdefined\HCode
        \HCode{%
<section class="learning-box learning-box--goals">
<div class="learning-box__header">
<span class="learning-box__icon" aria-hidden="true"></span>
<span class="learning-box__title">Leerdoelen</span>
</div>
<div class="learning-box__content">
        }%
        \begin{itemize}
    \else
\begin{tcolorbox}[
    enhanced,
    breakable,
    colback=fvdmagentalight,
    colframe=fvdmagenta,
    coltext=fvdtext,
    boxrule=0.5pt,
    leftrule=4pt,
    arc=4mm,
    outer arc=4mm,
    left=7mm,
    right=6mm,
    top=5mm,
    bottom=5mm,
    before skip=6mm,
    after skip=6mm,
    title=Leerdoelen,
    coltitle=fvdmagenta,
    fonttitle=\bfseries\Large,
    attach title to upper={\par\smallskip},
    boxed title style={
        empty,
        left=0mm,
        right=0mm,
        top=0mm,
        bottom=0mm
    },
    overlay={
        \draw[
            fvdmagenta,
            line width=0.5pt
        ]
        ([xshift=42mm,yshift=-13mm]frame.north west)
        --
        ([xshift=-7mm,yshift=-13mm]frame.north east);
    }
]
        \begin{itemize}
    \fi
}
{%
    \end{itemize}
    \ifdefined\HCode
        \HCode{%
</div>
</section>
        }%
    \else
        \end{tcolorbox}
    \fi
}


\addPrintStyle{../../..}

\title{Oefeningen --- Wiskundige uitspraken}
\author{Ingmar Herreman}

\begin{document}

% FVD_CHAPTER_DATA_START
% Automatisch gegenereerd uit index.tex.
% Niet handmatig aanpassen.
\fvdchapterdata{1}{Logica — De taal van correct redeneren}{1}
% FVD_CHAPTER_DATA_END













\maketitle


% ============================================================
% BASIS
% ============================================================


\begin{oefening}[Uitspraak, vraag of opdracht?]{\oefB \ESS}

Bepaal telkens of de gegeven zin een wiskundige uitspraak is.

\begin{question}
\[
13\text{ is een priemgetal.}
\]

\begin{multipleChoice}
  \choice[correct]{Dit is een wiskundige uitspraak.}
  \choice{Dit is geen wiskundige uitspraak.}
\end{multipleChoice}
\end{question}


\begin{question}
\[
\text{Bereken }7^2-3.
\]

\begin{multipleChoice}
  \choice{Dit is een wiskundige uitspraak.}
  \choice[correct]{Dit is geen wiskundige uitspraak.}
\end{multipleChoice}
\end{question}


\begin{question}
\[
\text{Is }21\text{ deelbaar door }7?
\]

\begin{multipleChoice}
  \choice{Dit is een wiskundige uitspraak.}
  \choice[correct]{Dit is geen wiskundige uitspraak.}
\end{multipleChoice}
\end{question}


\begin{question}
\[
21\text{ is deelbaar door }7.
\]

\begin{multipleChoice}
  \choice[correct]{Dit is een wiskundige uitspraak.}
  \choice{Dit is geen wiskundige uitspraak.}
\end{multipleChoice}
\end{question}


\begin{question}
Leg kort uit waarom de derde en vierde zin logisch niet van hetzelfde
type zijn.

\begin{oplossing}
De derde zin is een vraag. Hij vraagt naar een waarheidswaarde maar
beweert zelf niets.

De vierde zin is wel een bewering waarvan we kunnen bepalen of ze waar
of onwaar is. Het is dus een wiskundige uitspraak.
\end{oplossing}
\end{question}

\end{oefening}



\begin{oefening}[Waar of onwaar?]{\oefB \ESS}

Bepaal de waarheidswaarde van elke uitspraak en motiveer indien nodig.

\begin{question}
\[
3^3=27.
\]

\begin{oplossing}
De uitspraak is waar, want
\[
3^3=27.
\]
\end{oplossing}
\end{question}


\begin{question}
\[
\sqrt{49}=-7.
\]

\begin{oplossing}
De uitspraak is onwaar.

Met $\sqrt{49}$ bedoelen we de niet-negatieve vierkantswortel, dus
\[
\sqrt{49}=7.
\]
\end{oplossing}
\end{question}


\begin{question}
\[
17\in\mathbb{Q}.
\]

\begin{oplossing}
De uitspraak is waar. Elk geheel getal is rationaal, want
\[
17=\frac{17}{1}.
\]
\end{oplossing}
\end{question}


\begin{question}
\[
\pi\in\mathbb{Q}.
\]

\begin{oplossing}
De uitspraak is onwaar. $\pi$ is irrationaal.
\end{oplossing}
\end{question}


\begin{question}
\[
0\text{ is een positief reëel getal.}
\]

\begin{oplossing}
De uitspraak is onwaar. Het getal $0$ is noch positief, noch negatief.
\end{oplossing}
\end{question}

\end{oefening}



\begin{oefening}[Onwaar betekent niet ``geen uitspraak'']{\oefB \ESS}

Een leerling zegt:

\begin{quote}
``$15$ is een priemgetal is geen uitspraak, want het is fout.''
\end{quote}

Beoordeel deze redenering.

\begin{oplossing}
De redenering is fout.

De zin
\[
15\text{ is een priemgetal}
\]
is wel degelijk een wiskundige uitspraak, omdat we ondubbelzinnig
kunnen bepalen of hij waar of onwaar is.

De uitspraak is onwaar, omdat
\[
15=3\cdot5.
\]

Een uitspraak hoeft dus niet waar te zijn om een uitspraak te zijn.
\end{oplossing}

\end{oefening}



\begin{oefening}[Open uitspraak of concrete uitspraak?]{\oefB \ESS}

Bepaal telkens of het om een concrete uitspraak of een open uitspraak
gaat.

\begin{question}
\[
x+5=12.
\]

\begin{multipleChoice}
  \choice{Concrete uitspraak}
  \choice[correct]{Open uitspraak}
\end{multipleChoice}
\end{question}


\begin{question}
\[
7+5=12.
\]

\begin{multipleChoice}
  \choice[correct]{Concrete uitspraak}
  \choice{Open uitspraak}
\end{multipleChoice}
\end{question}


\begin{question}
\[
n\text{ is even.}
\]

\begin{multipleChoice}
  \choice{Concrete uitspraak}
  \choice[correct]{Open uitspraak}
\end{multipleChoice}
\end{question}


\begin{question}
\[
24\text{ is even.}
\]

\begin{multipleChoice}
  \choice[correct]{Concrete uitspraak}
  \choice{Open uitspraak}
\end{multipleChoice}
\end{question}


\begin{question}
Wat is het wezenlijke verschil tussen beide soorten?

\begin{oplossing}
Een concrete uitspraak heeft één vaste waarheidswaarde.

Bij een open uitspraak komt minstens één vrije veranderlijke voor en
hangt de waarheidswaarde af van de waarde die aan die veranderlijke
wordt toegekend.
\end{oplossing}
\end{question}

\end{oefening}



\begin{oefening}[Een predicaat onderzoeken]{\oefB \ESS}

Gegeven is

\[
P(x):\quad x^2<20.
\]

Bepaal telkens de waarheidswaarde.

\begin{question}
$P(2)$

\begin{oplossing}
\[
2^2=4<20.
\]

Dus $P(2)$ is waar.
\end{oplossing}
\end{question}


\begin{question}
$P(4)$

\begin{oplossing}
\[
4^2=16<20.
\]

Dus $P(4)$ is waar.
\end{oplossing}
\end{question}


\begin{question}
$P(5)$

\begin{oplossing}
\[
5^2=25>20.
\]

Dus $P(5)$ is onwaar.
\end{oplossing}
\end{question}


\begin{question}
$P(-4)$

\begin{oplossing}
\[
(-4)^2=16<20.
\]

Dus $P(-4)$ is waar.
\end{oplossing}
\end{question}


\begin{question}
Geef zelf één waarde waarvoor $P(x)$ waar is en één waarde waarvoor
$P(x)$ onwaar is.

\begin{oplossing}
Er zijn meerdere mogelijkheden.

Bijvoorbeeld:
\[
P(0)\text{ is waar}
\]
en
\[
P(10)\text{ is onwaar}.
\]
\end{oplossing}
\end{question}

\end{oefening}



\begin{oefening}[Zelf een predicaat formuleren]{\oefB \ESS}

Schrijf elke eigenschap als een predicaat.

\begin{question}
``$x$ is groter dan $7$.''

\begin{oplossing}
Bijvoorbeeld
\[
P(x):\quad x>7.
\]
\end{oplossing}
\end{question}


\begin{question}
``$n$ is deelbaar door $4$.''

\begin{oplossing}
Bijvoorbeeld
\[
Q(n):\quad n\text{ is deelbaar door }4.
\]
\end{oplossing}
\end{question}


\begin{question}
``Het kwadraat van $x$ is kleiner dan $10$.''

\begin{oplossing}
Bijvoorbeeld
\[
R(x):\quad x^2<10.
\]
\end{oplossing}
\end{question}


\begin{question}
Formuleer zelf een predicaat $S(n)$ over gehele getallen dat voor
sommige waarden waar en voor andere waarden onwaar is.

\begin{oplossing}
Er zijn veel mogelijkheden.

Bijvoorbeeld
\[
S(n):\quad n\text{ is oneven}.
\]

Dan is $S(5)$ waar en $S(8)$ onwaar.
\end{oplossing}
\end{question}

\end{oefening}



\begin{oefening}[Het domein verandert de situatie]{\oefB \ESS}

Beschouw

\[
P(x):\quad x^2=9.
\]

\begin{question}
Welke waarden voldoen als het domein $\mathbb{N}$ is?

\begin{oplossing}
In $\mathbb{N}$ voldoet
\[
x=3.
\]
\end{oplossing}
\end{question}


\begin{question}
Welke waarden voldoen als het domein $\mathbb{Z}$ is?

\begin{oplossing}
In $\mathbb{Z}$ voldoen
\[
x=-3
\qquad\text{en}\qquad
x=3.
\]
\end{oplossing}
\end{question}


\begin{question}
Wat leren deze twee antwoorden over het domein?

\begin{oplossing}
Het domein bepaalt welke waarden voor de veranderlijke toegelaten zijn.
Daardoor kan hetzelfde predicaat binnen verschillende domeinen tot
verschillende verzamelingen oplossingen leiden.
\end{oplossing}
\end{question}

\end{oefening}



\begin{oefening}[Kan de uitspraak waar worden?]{\oefB}

Beschouw de volgende predicaten.

\begin{question}
\[
P(x):\quad x^2=16,
\qquad x\in\mathbb{Z}.
\]

Geef alle waarden van $x$ waarvoor $P(x)$ waar is.

\begin{oplossing}
\[
x=-4\quad\text{of}\quad x=4.
\]
\end{oplossing}
\end{question}


\begin{question}
\[
Q(x):\quad x^2=3,
\qquad x\in\mathbb{Z}.
\]

Bestaat er een waarde waarvoor $Q(x)$ waar is?

\begin{oplossing}
Nee. Er bestaat geen geheel getal waarvan het kwadraat gelijk is aan
$3$.
\end{oplossing}
\end{question}


\begin{question}
Verander het domein van $Q(x)$ naar $\mathbb{R}$. Wat verandert er?

\begin{oplossing}
Nu zijn er wel waarden waarvoor het predicaat waar is:
\[
x=-\sqrt{3}
\qquad\text{en}\qquad
x=\sqrt{3}.
\]
\end{oplossing}
\end{question}

\end{oefening}



% ============================================================
% VERDIEPING
% ============================================================


\begin{oefening}[De formulering onderzoeken]{\oefU \ESS}

Een leerling schrijft:

\begin{center}
``$100$ is een groot getal.''
\end{center}

\begin{question}
Waarom is deze zin zonder extra informatie problematisch als
wiskundige uitspraak?

\begin{oplossing}
Het woord ``groot'' heeft zonder verdere afspraak geen ondubbelzinnige
wiskundige betekenis. De waarheidswaarde hangt af van de gekozen
context of norm.
\end{oplossing}
\end{question}


\begin{question}
Geef een preciezere wiskundige formulering waarin $100$ wel in een
ondubbelzinnige uitspraak voorkomt.

\begin{oplossing}
Bijvoorbeeld
\[
100>75.
\]

Of:
\[
100\text{ is deelbaar door }25.
\]

Andere ondubbelzinnige formuleringen zijn mogelijk.
\end{oplossing}
\end{question}

\end{oefening}



\begin{oefening}[Het verborgen domein]{\oefU \ESS}

Een leerling schrijft:

\begin{quote}
``De vergelijking $x^2=2$ heeft geen oplossing.''
\end{quote}

\begin{question}
Waarom is deze uitspraak onvoldoende precies?

\begin{oplossing}
Omdat niet vermeld wordt binnen welk domein naar oplossingen wordt
gezocht.

In $\mathbb{Q}$ heeft de vergelijking geen oplossing, maar in
$\mathbb{R}$ wel:
\[
x=\pm\sqrt{2}.
\]
\end{oplossing}
\end{question}


\begin{question}
Herschrijf de uitspraak zodat ze waar en volledig duidelijk is.

\begin{oplossing}
Bijvoorbeeld:
\[
\text{De vergelijking }x^2=2\text{ heeft geen oplossing in }\mathbb{Q}.
\]
\end{oplossing}
\end{question}


\begin{question}
Herschrijf ze vervolgens als een ware uitspraak over $\mathbb{R}$.

\begin{oplossing}
Bijvoorbeeld:
\[
\text{De vergelijking }x^2=2\text{ heeft twee oplossingen in }
\mathbb{R}.
\]

Die oplossingen zijn
\[
-\sqrt{2}\quad\text{en}\quad\sqrt{2}.
\]
\end{oplossing}
\end{question}

\end{oefening}



\begin{oefening}[Een fout in de redenering]{\oefU \ESS}

Een leerling onderzoekt de uitspraak

\begin{center}
``Het kwadraat van elk positief geheel getal is groter dan het getal
zelf.''
\end{center}

Hij berekent

\[
2^2>2,\qquad
3^2>3,\qquad
4^2>4,\qquad
5^2>5
\]

en besluit:

\begin{quote}
``Ik heb vier voorbeelden gecontroleerd, dus de uitspraak is bewezen.''
\end{quote}

\begin{question}
Wat is fout aan deze redenering?

\begin{oplossing}
Een eindig aantal bevestigende voorbeelden toont niet aan dat een
algemene uitspraak voor alle toegelaten waarden geldt.
\end{oplossing}
\end{question}


\begin{question}
Is de oorspronkelijke uitspraak waar?

\begin{oplossing}
Nee.

Voor
\[
n=1
\]
krijgen we
\[
1^2>1,
\]
dus
\[
1>1,
\]
wat onwaar is.

Het getal $1$ is een tegenvoorbeeld.
\end{oplossing}
\end{question}


\begin{question}
Wat zou veranderen als de bewering alleen betrekking had op gehele
getallen $n\geq2$?

\begin{oplossing}
Dan verdwijnt het tegenvoorbeeld $n=1$.

De uitspraak is voor $n\geq2$ inderdaad waar, maar de gecontroleerde
voorbeelden alleen vormen nog steeds geen algemeen bewijs.
\end{oplossing}
\end{question}

\end{oefening}



\begin{oefening}[Tegenvoorbeelden zoeken]{\oefU \ESS}

Geef indien mogelijk een tegenvoorbeeld voor elke bewering.
Als er geen tegenvoorbeeld bestaat, vermeld dat.

\begin{question}
Elk even natuurlijk getal is deelbaar door $4$.

\begin{oplossing}
Bijvoorbeeld
\[
2.
\]

Het getal $2$ is even maar niet deelbaar door $4$.
\end{oplossing}
\end{question}


\begin{question}
Elk veelvoud van $6$ is even.

\begin{oplossing}
Er bestaat geen tegenvoorbeeld.

Elk veelvoud van $6$ heeft de vorm
\[
6k=2(3k),
\]
en is dus even.
\end{oplossing}
\end{question}


\begin{question}
Voor elk reëel getal $x$ geldt
\[
x^2>x.
\]

\begin{oplossing}
Er zijn veel tegenvoorbeelden.

Bijvoorbeeld
\[
x=\frac12.
\]

Dan is
\[
x^2=\frac14<\frac12.
\]

Ook $x=0$ en $x=1$ zijn tegenvoorbeelden.
\end{oplossing}
\end{question}


\begin{question}
Elk priemgetal is oneven.

\begin{oplossing}
\[
2
\]
is een tegenvoorbeeld. Het is priem en even.
\end{oplossing}
\end{question}

\end{oefening}



\begin{oefening}[Een bewering herstellen]{\oefU \ESS}

Een algemene bewering blijkt fout te zijn.

Probeer ze telkens met een zo klein mogelijke wijziging correct te maken.

\begin{question}
Elk geheel getal is positief.

\begin{oplossing}
Er zijn verschillende mogelijke correcties.

Bijvoorbeeld:
\[
\text{Elk positief geheel getal is groter dan }0.
\]
\end{oplossing}
\end{question}


\begin{question}
Elk priemgetal is oneven.

\begin{oplossing}
Bijvoorbeeld:
\[
\text{Elk priemgetal groter dan }2\text{ is oneven.}
\]
\end{oplossing}
\end{question}


\begin{question}
Voor elk reëel getal $x$ geldt
\[
x^2>0.
\]

\begin{oplossing}
Bijvoorbeeld:
\[
\text{Voor elk reëel getal }x\neq0\text{ geldt }x^2>0.
\]

Een andere correcte formulering is
\[
x^2\geq0
\]
voor elk reëel getal $x$.
\end{oplossing}
\end{question}

\end{oefening}



\begin{oefening}[Dezelfde formule, andere betekenis]{\oefU \ESS}

Beschouw

\[
P(x):\quad x^2<4.
\]

Onderzoek het predicaat in verschillende domeinen.

\begin{question}
Voor welke $x\in\mathbb{Z}$ is $P(x)$ waar?

\begin{oplossing}
\[
x\in\{-1,0,1\}.
\]
\end{oplossing}
\end{question}


\begin{question}
Voor welke $x\in\mathbb{R}$ is $P(x)$ waar?

\begin{oplossing}
\[
-2<x<2.
\]
\end{oplossing}
\end{question}


\begin{question}
Leg uit waarom de formule zelf niet veranderd is, maar het antwoord
toch sterk verschilt.

\begin{oplossing}
Het predicaat $x^2<4$ blijft hetzelfde.

Het domein bepaalt echter welke waarden van $x$ toegelaten zijn.
Daarom is de verzameling waarden waarvoor het predicaat waar is
afhankelijk van het domein.
\end{oplossing}
\end{question}

\end{oefening}



\begin{oefening}[Wat wordt werkelijk beweerd?]{\oefU \ESS}

Bekijk de volgende drie zinnen:

\[
\begin{array}{ll}
\text{A.}&x^2=25,\\[1mm]
\text{B.}&5^2=25,\\[1mm]
\text{C.}&\text{Er bestaat een reëel getal waarvan het kwadraat }25\text{ is.}
\end{array}
\]

\begin{question}
Welke zin is een open uitspraak?

\begin{oplossing}
A is een open uitspraak, omdat $x$ vrij voorkomt.
\end{oplossing}
\end{question}


\begin{question}
Welke zinnen zijn concrete uitspraken met een waarheidswaarde?

\begin{oplossing}
B en C zijn concrete uitspraken.

Beide zijn waar.
\end{oplossing}
\end{question}


\begin{question}
Waarom is C logisch sterker geformuleerd dan enkel het invullen van
$x=5$?

\begin{oplossing}
C beweert dat er minstens één reëel getal bestaat dat aan de eigenschap
voldoet.

Het invullen van $x=5$ levert meteen zo'n voorbeeld en toont daardoor
dat C waar is.

Deze formulering loopt vooruit op de existentiële kwantor die later
wordt ingevoerd.
\end{oplossing}
\end{question}

\end{oefening}



\begin{oefening}[Een regelsysteem]{\oefU}

Een automatische ventilator in een serre gebruikt de voorwaarde

\[
P(T):\quad T>28,
\]

waarbij $T$ de temperatuur in graden Celsius is.

De ventilator wordt ingeschakeld wanneer $P(T)$ waar is.

\begin{question}
Bepaal de toestand van de ventilator bij
\[
T=24.
\]

\begin{oplossing}
Omdat
\[
24>28
\]
onwaar is, is $P(24)$ onwaar. De ventilator wordt niet ingeschakeld.
\end{oplossing}
\end{question}


\begin{question}
Bepaal de toestand bij
\[
T=31{,}5.
\]

\begin{oplossing}
Omdat
\[
31{,}5>28
\]
waar is, is $P(31{,}5)$ waar. De ventilator wordt ingeschakeld.
\end{oplossing}
\end{question}


\begin{question}
Een technicus schrijft in het programma:

\begin{center}
``Schakel de ventilator in als het warm is.''
\end{center}

Waarom is deze instructie ongeschikt voor een computerprogramma?

\begin{oplossing}
Het begrip ``warm'' is niet ondubbelzinnig gedefinieerd.
Een computerprogramma heeft een precieze voorwaarde nodig die als
waar of onwaar kan worden beoordeeld, bijvoorbeeld
\[
T>28.
\]
\end{oplossing}
\end{question}


\begin{question}
Welke rol speelt het domein van $T$ in dit model?

\begin{oplossing}
Het domein bepaalt welke temperaturen als mogelijke invoer worden
beschouwd. In een realistische toepassing kan men bijvoorbeeld een
praktisch meetbereik vastleggen.

De wiskundige voorwaarde krijgt betekenis binnen dat gekozen domein.
\end{oplossing}
\end{question}

\end{oefening}



\begin{oefening}[Wetenschappelijke uitspraken preciseren]{\oefU}

Beoordeel de volgende uitspraken als wetenschappelijke formulering.

\begin{question}
``Een zware planeet heeft een grote zwaartekracht.''

Waarom is deze zin onvoldoende precies om als wiskundige uitspraak te
gebruiken?

\begin{oplossing}
Begrippen zoals ``zwaar'' en ``groot'' zijn niet exact gedefinieerd.
Bovendien hangt de zwaartekracht niet alleen van de massa af, maar ook
van bijvoorbeeld de afstand tot het massamiddelpunt.

Er moet dus worden vastgelegd welke grootheden bedoeld worden en welke
voorwaarden gelden.
\end{oplossing}
\end{question}


\begin{question}
``De elektrische stroom is hoog.''

Welke informatie ontbreekt om hiervan een bruikbare logische
voorwaarde in een beveiligingssysteem te maken?

\begin{oplossing}
Er moet minstens een grenswaarde worden vastgelegd.

Bijvoorbeeld:
\[
I>10\,\mathrm{A}.
\]

Dan kan voor elke gemeten stroomsterkte ondubbelzinnig worden bepaald
of de voorwaarde waar of onwaar is.
\end{oplossing}
\end{question}

\end{oefening}



% ============================================================
% GEVORDERD
% ============================================================


\begin{oefening}[Een misleidend patroon]{\oefG}

Een leerling onderzoekt de formule

\[
P(n):\quad n^2+n+41\text{ is een priemgetal},
\qquad n\in\mathbb{N}.
\]

Voor

\[
n=0,1,2,3,4,5
\]

vindt hij telkens een priemgetal.

Daarom beweert hij:

\begin{quote}
``Voor elk natuurlijk getal $n$ is $n^2+n+41$ een priemgetal.''
\end{quote}

\begin{question}
Waarom vormen de zes berekeningen geen bewijs?

\begin{oplossing}
De bewering gaat over alle natuurlijke getallen. Het controleren van
een eindig aantal waarden bewijst niet dat het predicaat voor elke
mogelijke waarde waar is.
\end{oplossing}
\end{question}


\begin{question}
Onderzoek $n=40$.

\begin{oplossing}
Voor $n=40$ krijgen we
\[
40^2+40+41
=
1600+40+41
=
1681.
\]

Maar
\[
1681=41^2.
\]

Dit is dus geen priemgetal.
\end{oplossing}
\end{question}


\begin{question}
Wat is de logische betekenis van $n=40$ voor de oorspronkelijke
bewering?

\begin{oplossing}
$n=40$ is een tegenvoorbeeld.

Eén tegenvoorbeeld volstaat om de algemene bewering te weerleggen.
\end{oplossing}
\end{question}


\begin{question}
Welke belangrijke les over patroonherkenning volgt hieruit?

\begin{oplossing}
Een patroon dat voor veel voorbeelden zichtbaar is, kan aanleiding
geven tot een vermoeden maar vormt nog geen algemeen bewijs.

Bij algemene uitspraken moeten ook niet-geteste gevallen logisch worden
verantwoord.
\end{oplossing}
\end{question}

\end{oefening}



\begin{oefening}[Ontwerp zelf een tegenvoorbeeld]{\oefG}

Construeer zelf een foutieve algemene uitspraak over gehele getallen
die:

\begin{enumerate}
  \item voor minstens vijf eenvoudige testwaarden waar lijkt;
  \item toch minstens één tegenvoorbeeld heeft.
\end{enumerate}

Geef vervolgens:

\begin{enumerate}
  \item het gebruikte predicaat;
  \item het domein;
  \item vijf bevestigende voorbeelden;
  \item één tegenvoorbeeld;
  \item een verbeterde versie van de uitspraak.
\end{enumerate}

\begin{oplossing}
Er zijn veel correcte antwoorden mogelijk.

Een mogelijk voorbeeld is:

\[
P(n):\quad n^2+n+2\text{ is even},
\]
maar deze uitspraak is eigenlijk altijd waar en heeft dus geen
tegenvoorbeeld. Dat voorbeeld voldoet daarom niet aan de opdracht.

Een wel bruikbaar voorbeeld is:

\begin{quote}
Voor elk natuurlijk getal $n$ is $n^2+n+41$ priem.
\end{quote}

Voor verschillende kleine waarden van $n$ levert dit priemgetallen,
terwijl $n=40$ een tegenvoorbeeld geeft.

Een sterk antwoord vermeldt duidelijk waarom de eerste voorbeelden
alleen ondersteuning geven en waarom het tegenvoorbeeld de algemene
bewering ontkracht.
\end{oplossing}

\end{oefening}



\begin{oefening}[Een mini-onderzoek in logische precisie]{\oefG}

Kies één van de volgende contexten:

\begin{itemize}
  \item temperatuurregeling;
  \item batterijbewaking;
  \item verkeersdetectie;
  \item medische meting;
  \item robotica;
  \item een eigen wetenschappelijke of technologische context.
\end{itemize}

Formuleer binnen die context een predicaat

\[
P(x)
\]

dat door een automatisch systeem als waar of onwaar kan worden
beoordeeld.

Werk vervolgens de volgende punten uit:

\begin{enumerate}
  \item Wat stelt de veranderlijke $x$ voor?
  \item Wat is een zinvol domein?
  \item Wat betekent $P(x)$?
  \item Geef twee waarden waarvoor $P(x)$ waar is.
  \item Geef twee waarden waarvoor $P(x)$ onwaar is.
  \item Formuleer dezelfde voorwaarde eerst in gewone, maar vage taal.
  \item Leg uit waarom de wiskundige formulering beter geschikt is voor
        een automatisch systeem.
\end{enumerate}

\begin{oplossing}
Er zijn meerdere correcte oplossingen mogelijk.

Een mogelijk voorbeeld is batterijbewaking:

\[
P(V):\quad V<3{,}3,
\]

waarbij $V$ de batterijspanning in volt is.

Een mogelijk domein is bijvoorbeeld

\[
V\in[0,4{,}2].
\]

$P(3{,}0)$ en $P(2{,}8)$ zijn waar, terwijl $P(3{,}7)$ en $P(4{,}0)$
onwaar zijn.

De vage formulering ``de batterijspanning is laag'' wordt vervangen
door een precieze grensvoorwaarde. Daardoor kan het systeem voor elke
gemeten waarde ondubbelzinnig beslissen of de toestand optreedt.
\end{oplossing}

\end{oefening}



% ============================================================
% SYNTHESE
% ============================================================


\begin{oefening}[Alles samen: analyseer de bewering]{\oefU \ESS}

Beschouw de bewering:

\begin{center}
``Het kwadraat van elk geheel getal is groter dan het oorspronkelijke
getal.''
\end{center}

\begin{question}
Welke veranderlijke zou je gebruiken en wat is het domein?

\begin{oplossing}
Bijvoorbeeld $n$, met
\[
n\in\mathbb{Z}.
\]
\end{oplossing}
\end{question}


\begin{question}
Schrijf de bijhorende open uitspraak als predicaat.

\begin{oplossing}
\[
P(n):\quad n^2>n.
\]
\end{oplossing}
\end{question}


\begin{question}
Onderzoek $P(-2)$, $P(0)$, $P(1)$ en $P(3)$.

\begin{oplossing}
\[
P(-2):\quad4>-2
\]
is waar.

\[
P(0):\quad0>0
\]
is onwaar.

\[
P(1):\quad1>1
\]
is onwaar.

\[
P(3):\quad9>3
\]
is waar.
\end{oplossing}
\end{question}


\begin{question}
Is de oorspronkelijke algemene bewering waar?

\begin{oplossing}
Nee.

Zowel $n=0$ als $n=1$ is een tegenvoorbeeld.
\end{oplossing}
\end{question}


\begin{question}
Formuleer een aangepaste algemene uitspraak die wel waar is.

\begin{oplossing}
Bijvoorbeeld:

\begin{center}
Voor elk geheel getal $n<0$ of $n>1$ geldt $n^2>n$.
\end{center}

Of:
\[
n^2\geq n
\]
voor elk geheel getal $n$.
\end{oplossing}
\end{question}


\begin{question}
Waarom is deze oefening meer dan alleen een berekening?

\begin{oplossing}
We moeten de formulering interpreteren, een domein herkennen, een
predicaat opstellen, verschillende gevallen onderzoeken, een
tegenvoorbeeld gebruiken en uiteindelijk een algemene bewering
beoordelen en verbeteren.
\end{oplossing}
\end{question}

\end{oefening}



\end{document}